% Generated by build/sync_qft_soln.py from committed sources only.
% Source repository: https://github.com/Luca-Yucheng-Jin/QFT-soln
% Source commit: 07d7fe95fd5cf5b26d38a16ee535a04a925277b2
\section{The Faddeev--Popov Determinant (PSI QFT II, Tutorial 9)}

\noindent\textit{Question prompts paraphrased from the official
\href{https://psi-online.perimeterinstitute.ca/courses/take/quantum-field-theory-ii-student/pdfs/20776413-module-8-tutorial}{Perimeter PSI Online sheet};
the equations and notation follow the source.}

\subsection{The Faddeev--Popov Determinant}

As a finite-dimensional model of an unfixed gauge zero mode, take
\begin{equation}
    D=\begin{pmatrix}1&0\\0&0\end{pmatrix},
    \qquad
    A=\begin{pmatrix}a\\b\end{pmatrix}.
\end{equation}
Then
\begin{equation}
    I=\int_{-\infty}^{\infty}dA\,e^{-A^TDA}
    =\int_{-\infty}^{\infty}da\,db\,e^{-a^2}
\end{equation}
is ill-defined because one factor is the volume of the real line.

For the general gauge-invariant integral
\begin{equation}
    I=\int\mathcal{D}A\,e^{iS(A)},
\end{equation}
assume that $A\mapsto A_g$ belongs to a gauge group $\mathcal{G}$ and satisfies
\begin{equation}
    \mathcal{D}A=\mathcal{D}A_g,
    \qquad
    S(A)=S(A_g).
\end{equation}
Define the Faddeev--Popov determinant by
\begin{equation}
    1=\Delta(A)\int\mathcal{D}g\,
        \delta\!\bigl[f(A_g)\bigr],
    \qquad
    \int\mathcal{D}g=\operatorname{Vol}(\mathcal{G}).
\end{equation}

\begin{enumerate}[label=\alph*)]
    \item Prove that $\Delta(A)$ is gauge invariant.
    \begin{solution}
        Under a gauge transformation
        \begin{equation}
            A \mapsto A_h, \quad (A_h)_g = A_{gh}, \quad \Delta(A) \mapsto\Delta(A_h),
        \end{equation}
        we have
        \begin{equation}
            1 = \Delta(A_h)\int \mathcal{D}g\delta[f(A_{gh})].
        \end{equation}
        Formally, the functional integral over the gauge group is defined as
        \begin{equation}
            \mathcal{D}g = \prod_xd\mu_H(g(x)),
        \end{equation}
        where $d\mu_H$ is Haar measure. It is the group theoretical analogue to the normal translational invariant measure, $dx = d(x+a)$. In group theory, the abelian operation $+$ is replaced by actions like right or left actions. In our present case, we take the Haar measure to be right invariant, namely
        \begin{equation}
            \mathcal{D}g = \mathcal{D}(gh).
        \end{equation}
        Therefore, we have
        \begin{equation}
            1 = \Delta(A_h)\int \mathcal{D}(gh)\delta[f(A_{gh})] = \Delta(A_h)\int \mathcal{D}g'\delta[f(A_{g'})] = \Delta(A)\int \mathcal{D}g\delta[f(A_{g})] \implies \Delta(A) =\Delta(A_h).
        \end{equation}

    \end{solution}

    \item Insert the identity above and show that
    \begin{equation}
        I=\operatorname{Vol}(\mathcal{G})
        \int\mathcal{D}A\,e^{iS(A)}\Delta(A)\delta[f(A)].
    \end{equation}
    \begin{solution}
        Inserting the identity into the path integral, we find
        \begin{equation}
            I = \int \mathcal{D}A e^{iS[A]}\Delta(A)\int\mathcal{D}g\delta[f(A_g)].
        \end{equation}
        Using gauge invariant of the action and the integral measure, $\mathcal{D}A = \mathcal{D}A_g$, we find
        \begin{equation}
            I = \left(\int \mathcal{D}g\right)\int\mathcal{D}A_g\,e^{iS(A_g)}\Delta(A_g)\delta[f(A_g)] = \operatorname{Vol}(\mathcal{G})
        \int\mathcal{D}A\,e^{iS(A)}\Delta(A)\delta[f(A)].
        \end{equation}
    \end{solution}
\end{enumerate}

Next use the zero-dimensional toy theory
\begin{equation}
    Z=\int_{-\infty}^{\infty}dx\,dy\,e^{-S(x,y)},
    \qquad
    S(x,y)=(x-y)^2.
\end{equation}

\begin{enumerate}[label=\alph*),resume]
    \item Identify the global ``gauge'' symmetry of this zero-dimensional action.
    \begin{solution}
        The toy theory is clearly invariant under translation, $x \to x+a, y \to y+a$.
    \end{solution}

    \item Describe the full configuration space $\mathcal{A}$ and its gauge
    orbits $\mathcal{G}$.
    \begin{solution}
        \begin{equation}
            \mathcal{A} = \{x,y\in \mathbb{R}\} = \mathbb{R}^2, \quad \mathcal{G} = (\mathbb{R},+) \implies \mathcal{A}/\mathcal{G} \cong \mathbb{R}.
        \end{equation}
    \end{solution}
\end{enumerate}

The divergent partition function can be factorized as
\begin{equation}
    Z=\left(\int da\right)X,
\end{equation}
where $a$ parametrizes the gauge group and $X$ is independent of $a$.

\begin{enumerate}[label=\alph*),resume]
    \item Find new variables that make this factorization explicit. As an
    additional challenge, suggest a modified model whose gauge-group volume is
    finite.
    \begin{solution}
        Perform the change of variable,
        \begin{equation}
            a = \frac{x-y}{\sqrt{2}}, \quad b = \frac{x+y}{\sqrt{2} },
        \end{equation}
        which gives
        \begin{equation}
            Z = \left(\int da\right)\int db e^{-b^2}.
        \end{equation}
    \end{solution}

    \item Choose a gauge slice for computing $X$ and give a condition that
    characterizes it; the choice need not be unique. Denote the condition by
    $f(x,y)=0$.
    \begin{solution}
        It is easiest if we choose the gauge slice $f(x,y) = y = 0$.
    \end{solution}

    \item Show that
    \begin{equation}
        Z=\int_{-\infty}^{\infty}da
        \int_{-\infty}^{\infty}dx\,dy\,
        J\,\delta\!\bigl(f(x,y)\bigr)e^{-(x-y)^2},
    \end{equation}
    where the Jacobian $J$ is independent of $a$.
    \begin{solution}
        We have
        \begin{equation}
            1 = \int df(x,y) \delta(f(x,y)) = \int df(x+a,y+a) \delta(f(x+a,y+a)) = \int da \, J\delta(f(x+a,y+a)).
        \end{equation}
        For our particular gauge choice, this is
        \begin{equation}
            1 = \int da \delta(y+a).
        \end{equation}
        Inserting into the partition function, we find
        \begin{equation}
            Z=\int_{-\infty}^{\infty}da
        \int_{-\infty}^{\infty}dx\,dy\,
        J\,\delta\!\bigl(f(x,y)\bigr)e^{-(x-y)^2}
        \end{equation}
        by "gauge" invariance of the action and the measure.
    \end{solution}
\end{enumerate}

Replace the gauge slice $f=0$ by $f=c$, with constant $c\neq0$, so that
\begin{equation}
    Z=\int da\int dx\,dy\,
    J\,\delta\!\bigl(f(x,y)-c\bigr)e^{-(x-y)^2}.
\end{equation}

\begin{enumerate}[label=\alph*),resume]
    \item Although the last expression displays $c$ explicitly, prove that $Z$
    is independent of $c$.
    \begin{solution}
        Since
        \begin{equation}
            1 = \int df(x,y) \delta(f(x,y)) = 1 = \int df(x,y) \delta(f(x,y)-c),
        \end{equation}
        the rest follows straightforwardly from above.
    \end{solution}

    \item Average over $c$ with the Gaussian weight
    \begin{equation}
        \int_{-\infty}^{\infty}dc\,e^{-c^2}
    \end{equation}
    and derive
    \begin{equation}
        Z=\frac{1}{\sqrt{\pi}}
        \int da\int dx\,dy\,
        J\,e^{-f^2}e^{-(x-y)^2}.
    \end{equation}
    Explain how the factor $e^{-f^2}$ supplies a gauge-fixing term.
    \begin{solution}
        \begin{equation}
            \begin{aligned}
                Z=\frac{1}{\sqrt{\pi}}\int dc\, e^{-c^2}\int da\int dx\,dy\,
    J\,\delta\!\bigl(f(x,y)-c\bigr)e^{-(x-y)^2} = \frac{1}{\sqrt{\pi}}
        \int da\int dx\,dy\,
        J\,e^{-f^2}e^{-(x-y)^2}.
            \end{aligned}
        \end{equation}
    This factor is largest when \(f=0\) and exponentially suppresses configurations for which \(f\neq0\). It therefore imposes the same gauge condition softly.
    \end{solution}

    \item Exponentiate the remaining Jacobian by representing it as an integral
    over Grassmann ``ghost'' variables.
    \begin{solution}
        \begin{equation}
            J = \int d\bar cdc \, e^{-\bar c J c}.
        \end{equation}
    \end{solution}
\end{enumerate}
